\chapter{完整源码清单：LCQI 本地 CPU 推理引擎}

本附录完整收录 \filepath{books/cpu-volume-3/labs/linux_cpu_inference} 的主线工程。第三册的代码主线是 LCQI：从模型 manifest、tokenizer、tiny MLP、int8 kernel、reference decoder trace、7B ledger 到 serving probe。读者必须把每个文件看成全链路的一段，而不是孤立 demo。

\section{一、CMakeLists.txt}

CMake 文件定义 LCQI 库、CLI、trace、ledger、serving probe、benchmark、可选 oneDNN 对标和测试目标。读者应从这里看哪些能力是必构建，哪些是平台可选。

\lstinputlisting[language=bash,caption={linux\_cpu\_inference/CMakeLists.txt}]{../../labs/linux_cpu_inference/CMakeLists.txt}

\section{二、README}

README 给出本地构建、测试、CLI、benchmark、trace、ledger、serving probe 和真实 small 模型 smoke 的命令。它是复现实验的入口，也说明 tiny 模型、真实 small 模型 smoke 和未来 7B 目标之间的边界。

\lstinputlisting[language=bash,caption={linux\_cpu\_inference/README.md}]{../../labs/linux_cpu_inference/README.md}

\section{三、真实模型 Smoke 工具}

\filepath{tools/run_smollm2_small_smoke.py} 下载并校验 SmolLM2-135M-Instruct 的 GGUF 量化文件，构建固定 commit 的 \code{llama-simple-chat}，在 CPU 上运行一次真实开源 small 模型生成，并输出可复查报告。它是外部真实模型 smoke，不替代 LCQI tiny golden trace。

\lstinputlisting[language=Python,caption={tools/run\_smollm2\_small\_smoke.py}]{../../tools/run_smollm2_small_smoke.py}

\filepath{tools/run_gpt2_smoke.py} 下载标准 GPT-2 的 \code{config.json}、\code{vocab.json}、\code{merges.txt} 和 \code{model.safetensors}，再调用 LCQI 自研 \code{lcqi_gpt2} reference path 生成一个 token。它证明的是自研 safetensors、BPE、GPT-2 forward 和 greedy generation 已连成闭环，不是高性能 benchmark。

\lstinputlisting[language=Python,caption={tools/run\_gpt2\_smoke.py}]{../../tools/run_gpt2_smoke.py}

\section{四、模型公共头}

\filepath{model.hpp} 定义 tiny MLP 模型结构、权重、bias、scale 和加载入口。它负责说明模型文件被解释成哪些 C++ 对象。

\lstinputlisting[language=C++,caption={include/lcqi/model.hpp}]{../../labs/linux_cpu_inference/include/lcqi/model.hpp}

\section{五、推理公共头}

\filepath{inference.hpp} 定义推理结果和 \code{run_inference} 入口。读者要看清这里返回 hidden、logits 和 predicted class，因为测试会用这些字段定位闭环错误。

\lstinputlisting[language=C++,caption={include/lcqi/inference.hpp}]{../../labs/linux_cpu_inference/include/lcqi/inference.hpp}

\section{六、Reference Decoder 公共头}

reference decoder 头文件定义张量、层配置、trace checkpoint、采样结果和 decode 函数。它是定位推理错误的主要合同，不只返回最终 token。

\lstinputlisting[language=C++,caption={include/lcqi/reference\_decoder.hpp}]{../../labs/linux_cpu_inference/include/lcqi/reference_decoder.hpp}

\section{七、GPT-2 Reference 公共头}

GPT-2 reference 头文件定义 GPT-2 配置、byte-level BPE tokenizer、HuggingFace safetensors 加载、forward 和 greedy generation。它单独存在，是因为 GPT-2 使用 LayerNorm、绝对位置嵌入、packed QKV、GELU 和 tied lm head，不能直接塞进 LLaMA-like reference decoder。

\lstinputlisting[language=C++,caption={include/lcqi/gpt2\_reference.hpp}]{../../labs/linux_cpu_inference/include/lcqi/gpt2_reference.hpp}

\section{八、Int8 Kernel 公共头}

int8 kernel 头文件定义 packed linear、deterministic fixture、scalar/workspace/AVX2 路径和误差比较。读者应注意 AVX2 可用性是运行时事实，不能只看源码存在。

\lstinputlisting[language=C++,caption={include/lcqi/int8\_kernels.hpp}]{../../labs/linux_cpu_inference/include/lcqi/int8_kernels.hpp}

\section{九、Tokenizer 公共头}

tokenizer 头文件固定 BOS/EOS/UNK、词表和 template hash。读者应重点看未知 token 行为是否明确，因为 prompt 合同不稳定会污染后续所有 trace。

\lstinputlisting[language=C++,caption={include/lcqi/tokenizer.hpp}]{../../labs/linux_cpu_inference/include/lcqi/tokenizer.hpp}

\section{十、Safetensors 公共头}

safetensors 头文件固定 dtype、shape、offset 和 payload 边界。它负责说明模型字节怎样被解释，不能和数学推理逻辑混在一起。

\lstinputlisting[language=C++,caption={include/lcqi/safetensors.hpp}]{../../labs/linux_cpu_inference/include/lcqi/safetensors.hpp}

\section{十一、模型加载实现}

模型加载实现负责解析 tiny 文本模型、权重、scale、bias 和 shape。读者应检查坏格式、维度和数值解析是否有清晰错误边界。

\lstinputlisting[language=C++,caption={src/model.cpp}]{../../labs/linux_cpu_inference/src/model.cpp}

\section{十二、Tiny MLP 推理实现}

推理实现执行 int8 linear、ReLU、第二层 logits 和 predicted class。它是最短闭环，不是完整 LLM；价值在于把模型加载和 kernel 调用接起来。

\lstinputlisting[language=C++,caption={src/inference.cpp}]{../../labs/linux_cpu_inference/src/inference.cpp}

\section{十三、命令行入口}

\filepath{src/main.cpp} 把模型路径和输入接到 \code{run_inference}，并打印结果。CLI 应保持薄层，不复制模型解析和推理逻辑。

\lstinputlisting[language=C++,caption={src/main.cpp}]{../../labs/linux_cpu_inference/src/main.cpp}

\filepath{src/gpt2_cli.cpp} 负责把模型目录、prompt 和 \code{max\_new\_tokens} 接到 GPT-2 reference path。无参数时运行 tiny GPT-2 fixture；传入 HuggingFace 模型目录时加载真实资产。

\lstinputlisting[language=C++,caption={src/gpt2\_cli.cpp}]{../../labs/linux_cpu_inference/src/gpt2_cli.cpp}

\section{十四、Reference Decoder 完整实现}

这是第三册最重要的长文件。它实现 embedding、RMSNorm、Q/K/V、RoPE、KV cache、attention、SwiGLU、final norm、logits 和 sampler，并在关键位置生成 trace checkpoint。读者应按 checkpoint 顺序读，而不是从最终 logits 倒推。

\lstinputlisting[language=C++,caption={src/reference\_decoder.cpp}]{../../labs/linux_cpu_inference/src/reference_decoder.cpp}

\section{十五、GPT-2 Reference 完整实现}

GPT-2 reference 实现包括 JSON 配置读取、byte-level BPE、safetensors F32/F16/BF16 张量读取、两套 HuggingFace 权重名前缀解析、LayerNorm、packed QKV、causal attention、GELU、tied lm head 和 greedy generation。它是正确性优先的 reference path，暂不追求吞吐。

\lstinputlisting[language=C++,caption={src/gpt2\_reference.cpp}]{../../labs/linux_cpu_inference/src/gpt2_reference.cpp}

\section{十六、Int8 Kernel 基础实现}

\filepath{int8_kernels.cpp} 包含 scalar、packed、workspace、deterministic fixture 和误差比较。它是 AVX2 路径的 reference 对照。

\lstinputlisting[language=C++,caption={src/int8\_kernels.cpp}]{../../labs/linux_cpu_inference/src/int8_kernels.cpp}

\section{十七、Int8 Kernel AVX2 实现}

\filepath{int8_kernels_avx2.cpp} 是平台优化路径。读者要先看可用性判断，再看向量化实现；如果机器不支持 AVX2，报告必须写 unavailable。

\lstinputlisting[language=C++,caption={src/int8\_kernels\_avx2.cpp}]{../../labs/linux_cpu_inference/src/int8_kernels_avx2.cpp}

\section{十八、Tokenizer 实现}

tokenizer 实现把 prompt 变成 token id，并固定 unknown token 行为。读者应检查 BOS/EOS 是否稳定插入，UNK 是否按合同处理。

\lstinputlisting[language=C++,caption={src/tokenizer.cpp}]{../../labs/linux_cpu_inference/src/tokenizer.cpp}

\section{十九、Safetensors 实现}

safetensors 实现解析 header、metadata、dtype、shape 和 data offsets。坏 dtype、坏 shape、坏 offset 和 payload 太短都应该被拒绝。

\lstinputlisting[language=C++,caption={src/safetensors.cpp}]{../../labs/linux_cpu_inference/src/safetensors.cpp}

\section{二十、Benchmark 程序}

\filepath{bench.cpp} 输出 int8 kernel benchmark。它服务性能观察，但必须和 max diff 对照一起看。

\lstinputlisting[language=C++,caption={src/bench.cpp}]{../../labs/linux_cpu_inference/src/bench.cpp}

\section{二十一、Trace 程序}

\filepath{trace.cpp} 输出 reference decoder checkpoint。它用于定位错误边界，不应该被简化成只打印最终 token。

\lstinputlisting[language=C++,caption={src/trace.cpp}]{../../labs/linux_cpu_inference/src/trace.cpp}

\section{二十二、Ledger 程序}

\filepath{ledger.cpp} 输出 7B 规模账本。它是估算，不是实测吞吐；读者要区分 FLOPs、KV 容量和带宽约束。

\lstinputlisting[language=C++,caption={src/ledger.cpp}]{../../labs/linux_cpu_inference/src/ledger.cpp}

\section{二十三、Serving Probe 程序}

\filepath{serving_probe.cpp} 固定短请求、RAG 请求、取消请求和超长请求。它把 TTFT、TPOT、拒绝和取消回收变成可观察字段。

\lstinputlisting[language=C++,caption={src/serving\_probe.cpp}]{../../labs/linux_cpu_inference/src/serving_probe.cpp}

\section{二十四、oneDNN 对标程序}

\filepath{onednn_bench.cpp} 是可选外部库对标入口。它只在环境安装 oneDNN 时构建，不能把不可用平台写成已验证性能。

\lstinputlisting[language=C++,caption={src/onednn\_bench.cpp}]{../../labs/linux_cpu_inference/src/onednn_bench.cpp}

\section{二十五、Tiny MLP Fixture}

\filepath{tiny_mlp_i8.txt} 是 tiny MLP 的模型资产。它决定 shape、权重、scale 和 bias，测试中的 logits golden 依赖它。

\lstinputlisting[language=bash,caption={models/tiny\_mlp\_i8.txt}]{../../labs/linux_cpu_inference/models/tiny_mlp_i8.txt}

\section{二十六、Tiny Tokenizer Fixture}

\filepath{tiny_tokenizer.txt} 决定词表、BOS/EOS/UNK 和模板 hash。prompt 变 token 的合同从这里开始。

\lstinputlisting[language=bash,caption={models/tiny\_tokenizer.txt}]{../../labs/linux_cpu_inference/models/tiny_tokenizer.txt}

\section{二十七、完整测试}

测试文件把 safetensors、tokenizer、tiny MLP、int8 kernel、reference decoder、trace、ledger 和 serving probe 串成验收网。新增优化或 shape 时，必须先扩展这里的 golden 和边界样例。

\lstinputlisting[language=C++,caption={tests/lcqi\_tests.cpp}]{../../labs/linux_cpu_inference/tests/lcqi_tests.cpp}
